\subsection{The exact denominator of a chain}
\label{shared-level-chain}

The auxiliary-prime calculation can keep each step's cost and the
positive factor already present in the source lattice.  The following
identity does not require the intermediate operators to be invertible.

\begin{proposition}\label{shared-level-chain-integrality}
Let $R$ be a discrete valuation ring with fraction field $K$, normalized
valuation $v$, and a field extension $E/K$.
Fix $q\in R\setminus\{0\}$, with $e=v(q)>0$.
For $0\le i\le m$, let $\Lambda_i$ be a finitely generated
$R$-submodule of a finite-dimensional $E$-space $V_i$, and let
$U_i\in\operatorname{End}_E(V_i)$.
Let $F_i:V_{i-1}\to V_i$ satisfy
\[
 U_iF_i=q^{r_i}F_iU_{i-1},\qquad
 A_i:=q^{d_i}U_iF_i,\qquad A_i\Lambda_{i-1}\subseteq\Lambda_i
 \quad(1\le i\le m),
\]
where $r_i,d_i\in\mathbb Z$.  If
$M\subseteq q^aU_0^m\Lambda_0$, with $a\in\mathbb Z$, put
\[
 F=F_m\cdots F_1,\qquad
 \ell=\sum_{i=1}^m d_i+
           \sum_{i=1}^m(m-i+1)r_i-a.
\]
Then $q^\ell F(M)\subseteq\Lambda_m$.  More generally,
$cF(M)\subseteq\Lambda_m$ whenever $c\in K^\times$ and
$v(c)\ge e\ell$.  Negative values of $\ell$ are allowed.
\end{proposition}
\begin{proof}
Repeatedly move an operator $U_i$ across the maps to its right.  Each
commutation with $F_i$ occurs $m-i+1$ times, giving
\[
 A_m\cdots A_1
 =q^{\sum_i d_i+\sum_i(m-i+1)r_i}F U_0^m.
\]
The left side takes $\Lambda_0$ into $\Lambda_m$.  For
$x=q^aU_0^my\in M$, this identity gives
$q^\ell F(x)=A_m\cdots A_1(y)\in\Lambda_m$.
Finally $v(cq^{-\ell})\ge0$, so multiplication by this element
preserves $\Lambda_m$.
\end{proof}

This exponent is optimal from the displayed assumptions.  In rank one,
take $\Lambda_i=R$, $U_0=1$,
$U_i=q^{r_1+\cdots+r_i}$, and $F_i=q^{-d_i}U_i^{-1}$.
Then $A_i=1$ and, for $M=q^aR$, one has
$F(M)=q^{-\ell}R$.  The empty chain also gives the asserted formula.

This formulation permits complex ambient function spaces with their
embedded integral lattices. Membership in the original rational image
supplies the vector $y$ used in the proof; each integral $A_i$ then
supplies the next rational lift. No separate rationality of $F_i$ or
$U_i$ is assumed.

For the arithmetic chain, let $T$ be the auxiliary rational prime,
$t$ odd, $h=(t+1)/2$, and $J$ the selected set of coefficient
embeddings, of size $j$.  At an active place above $T$, the prescribed
source lattice is contained in
\[
 T^a U_0^m\Lambda_0,\qquad
 a=m\sum_{\tau\in J}b_\tau.
\]
The original local kernel and value-lattice identities give the
single-step integers
\[
 d_i=fN_{i-1}+\sum_\tau s_{i-1,\tau}
                    +\sum_{\tau\notin J}\lambda_{i,\tau},
 \qquad r_i=hj.
\]
Here $s_{i-1,\tau}$ is the determinant shift to be cleared and
$\lambda_{i,\tau}$ is the first residual highest weight after that
shift.  These identities concern the actual coefficient lattices and
the actual kernels; the stabilizer orders in their finite sums must
be units at the place.  The resulting denominator is therefore
\begin{equation}\label{shared-level-chain-exact-loss}
 \ell=\sum_{i=1}^m d_i+hj\frac{m(m+1)}2
                         -m\sum_{\tau\in J}b_\tau.
\end{equation}
If the original weight bounds give $d_i\le2W+fN_{i-1}$ and
$N_i=k+i(t+1)$, then the former common-step upper bound
$B=m(2W+fN_m)+hj m(m+1)/2$ improves to
\[
 \ell\le B-\frac{f(t+1)m(m+1)}2-m\sum_{\tau\in J}b_\tau.
\]
The gain uses the same source lattice; its degree is not changed or
counted a second time.  At inactive places, the integral-step argument
still gives exponent zero.  The local kernel identities, rational
realizations, twisted-chain identifications and all other period
denominators remain required.  This computation supplies the passage
from those local identities to the full-chain bound, not their
arithmetic construction.
